% ============================================================
%  GENERAL INTRODUCTION
% ============================================================

\chapter*{General Introduction}
\label{ch:intoGen}
\addcontentsline{toc}{chapter}{General Introduction}
\markboth{General Introduction}{General Introduction}

Agriculture plays a central role in food security, rural livelihoods, and
economic planning. In Rwanda, reliable information on crop type distribution
is essential for agricultural monitoring, input planning, yield estimation,
early warning systems, and national food-security decision-making. However,
producing accurate crop maps remains challenging. Rwandan agricultural
landscapes are characterised by smallholder fields, fragmented land parcels,
mixed farming practices, regional differences in crop calendars, and frequent
cloud cover during the growing season. These conditions make crop type mapping
a complex remote-sensing problem, particularly when the goal is to distinguish
between individual crop classes rather than simply identify cropland.

Satellite remote sensing offers an effective way to observe agricultural areas
regularly and at large scale. Sentinel-2 imagery is especially useful for this
task because it provides multispectral observations at a spatial resolution
suitable for agricultural monitoring. Its visible, red-edge, near-infrared, and
shortwave-infrared bands capture information related to vegetation vigour,
canopy structure, moisture, and crop development. When analysed as a time
series, these observations describe how crops evolve across the growing season.
This temporal behaviour is important because different crops may look similar
on a single date but follow different seasonal trajectories.

This project investigates the use of Sentinel-2 time series and supervised
machine learning for multi-class crop type mapping in Rwanda. The study focuses
on Season~B~2025 and uses the Harvard Dataverse Rwanda crop mapping ground
reference dataset as the main source of labelled crop polygons. The
classification task focuses on four monocrop classes: maize, rice, bean, and
Irish potato. These crops are important in the Rwandan agricultural context and
provide a meaningful case study for evaluating whether satellite time-series
data can support operational crop monitoring.

The main aim of this project is to design, implement, and evaluate a
reproducible crop classification workflow that transforms field-labelled crop
polygons and Sentinel-2 observations into reliable crop type predictions. The
project is not limited to training a single model; it develops a complete
analytical pipeline that addresses the practical challenges of working with
real agricultural remote-sensing data. These include data quality issues, cloud
contamination, temporal gaps, class imbalance, spatial dependence, and model
comparison. This makes the work valuable both as a modelling study and as a
structured case study in building an Earth Observation pipeline for crop
monitoring.

The project has four main objectives. First, the ground reference and satellite
time-series data are prepared and quality-controlled to produce a consistent
modelling dataset. Second, the seasonal behaviour of each crop is represented
using spectral, vegetation-index, and phenological information. Third,
classical machine-learning models and sequence-based deep learning models are
compared under a consistent evaluation framework. Fourth, model performance is
assessed while accounting for class imbalance and possible spatial leakage, so
that the results provide a more realistic view of generalisation.

A key contribution of this work is its emphasis on methodological
reproducibility. Each stage of the workflow is documented, from data
preparation and temporal reconstruction to feature extraction, model
development, and evaluation. This transparency is important because, in remote
sensing, model performance depends strongly on how satellite observations and
ground labels are prepared before training begins. By making these decisions
explicit, the project provides a workflow that can be reviewed, adapted, and
extended in future crop monitoring studies.

The project also has practical relevance. Crop maps can support government
agencies, development organisations, and agricultural stakeholders by improving
seasonal monitoring, guiding agricultural planning, and strengthening
food-security assessments. At the same time, the project recognises that
machine-learning outputs must be interpreted carefully. Minority-class
performance, spatial generalisation, and single-season training data are
therefore treated as important methodological concerns rather than ignored
limitations.

The outcome of this project is a complete and reproducible framework for
multi-class crop type classification from Sentinel-2 time series in Rwanda. It
demonstrates how field-collected crop labels can be linked with satellite
observations, transformed into meaningful temporal representations, and
evaluated through a comparative modelling framework. More broadly, the work
contributes to the development of practical Earth Observation methods for
agricultural monitoring in smallholder farming systems, where high-quality crop
information is valuable but difficult to obtain at scale.

The remainder of this report is organised as follows.
Chapter~\ref{ch:introduction} introduces the host organisation, the project
background, the problem statement, the aims and objectives, and the project
management plan. Chapter~\ref{ch:literature} reviews the literature on remote
sensing, Sentinel-2 imagery, crop phenology, crop classification methods, and
class imbalance handling. Chapter~\ref{ch:method} presents the methodology,
including dataset preparation, preprocessing, feature engineering, model
development, and evaluation design. Chapter~\ref{ch:results} presents the
experimental results. Chapter~\ref{ch:discussion} discusses the significance
of the findings, the strengths and limitations of the work, and their
implications. The final chapter concludes the report and outlines possible
directions for future work.